\begin{table}[H]
\centering
\caption{Cancer groups with significant age decile survival differences
(multi-group log-rank test per group; BH-FDR q-value adjusted across cancer
groups within each phase; three significance tiers).\label{tab:age_decile_sig}}

\noindent\textbf{OS (12/22 significant):}
\begin{itemize}\setlength{\itemsep}{0pt}
  \item \textbf{q$ < $0.001:} High-grade glioma, Diffuse midline glioma, Choroid plexus tumor, Low-grade glioma, Medulloblastoma, Ewing sarcoma
  \item \textbf{q$ < $0.01:} Diffuse hemispheric glioma, Atypical Teratoid Rhabdoid Tumor
  \item \textbf{q$ < $0.05:} Glial-neuronal tumor NOS, Sarcoma, CNS Embryonal tumor, Chordoma
\end{itemize}

\noindent\textbf{EFS (17/22 significant):}
\begin{itemize}\setlength{\itemsep}{0pt}
  \item \textbf{q$ < $0.001:} Adamantinomatous Craniopharyngioma, Diffuse midline glioma, High-grade glioma, Low-grade glioma, Medulloblastoma, Meningioma, Ependymoma
  \item \textbf{q$ < $0.01:} Diffuse hemispheric glioma, Choroid plexus tumor, Embryonal tumor with multilayer rosettes, Schwannoma, Glial-neuronal tumor NOS
  \item \textbf{q$ < $0.05:} Chordoma, Atypical Teratoid Rhabdoid Tumor, Neurofibroma/Plexiform, Dysembryoplastic neuroepithelial tumor, Pineoblastoma
\end{itemize}
\end{table}
